\subsection{The remaining assertion about local ind-presentations}

\begin{claim}{E.17}{Open}{V ms.~44 and 46: local presentations of invariant ind-objects}
The following sequential assertion remains unresolved in this audit.
Let $K=(K_n)_{n\in\N}$ be a sequence in $\Cathot_\U$, and assume
that its associated presheaf
\[
 H_K(X)=\colim_n\Hom_{\Cathot}(X,K_n)
\]
sends every map in $\W_\infty$ between $\U$-small categories to a
bijection. Must $K$ be isomorphic in $\Ind_\U(\Cathot_\U)$ to a
sequence of local categories, equivalently to a sequence of groupoids?
The corresponding assertion with an arbitrary $\U$-small filtered
indexing category in place of $\N$ also remains unresolved here.
We do not infer either assertion from descent in A.7, from the
representable-object classification in Part C, or from the obstruction
in D.8 for a single represented homotopy type.
\end{claim}

There is, however, a proved exact criterion for both versions. Let
$\mathcal G\subset\Cathot_\U$ be the full subcategory on groupoids,
let $j:\mathcal G\hookrightarrow\Cathot_\U$ be its inclusion, and
let $\eta_X:X\to j\Pi_1X$ be the unit of the reflection in C.4.
For any $\U$-small filtered $I$ and $K:I\to\Cathot_\U$, the
following conditions are equivalent:
\[
\begin{gathered}
 H_K(\eta_X)\text{ is bijective for every }X\in\Cat_\U;\\
 K\longrightarrow(\Ind_\U j)(\Ind_\U\Pi_1)K
       \text{ is an isomorphism};\\
 K\text{ is isomorphic to a }\U\text{-small ind-object of groupoids}.
\end{gathered}
\tag{E.17a}
\]
These conditions are also equivalent to invariance under $\W_1$,
the functors inducing bijections on components and isomorphisms on
fundamental groups. When they hold, the groupoid presentation can
be taken to be the canonical diagram $(\Pi_1K_i)_{i\in I}$ on the
\emph{same} indexing category. In particular the sequential assertion
does not require an unproved change from a general filtered index
back to a sequential one.

\begin{proof}[Proof of the criterion and the verified special cases]
C.4 gives the fully faithful right adjoint $j$ to $\Pi_1$ and
identifies its inverted class with $\W_1$. Apply A.4 to this actual
adjunction. It gives (E.17a), its $\W_1$ formulation, and the
specified unit. A diagram of local categories is equivalent to a
diagram of groupoids by applying its objectwise invertible units,
using C.3. The canonical comparison on presheaves is explicitly
\[
 \colim_i\Hom_{\Cathot}(X,K_i)
 \longrightarrow
 \colim_i\Hom_{\mathcal G}(\Pi_1X,\Pi_1K_i).
 \tag{E.17b}
\]

If the initial hypothesis of E.17 holds and the $BK_i$ are
$1$-types on a cofinal subdiagram, A.7 identifies the left side
of (E.17b) with $\colim_i[BX,BK_i]$. Each $B\eta_{K_i}$ on
that subdiagram is a weak equivalence, and groupoid locality
identifies the right side with the same colimit. Thus (E.17b)
is bijective and the conclusion holds in this case. Essentially
constant ind-objects with invariant presheaves are also covered
directly by A.7 and C.3, after replacing the diagram by a constant
presentation. More generally, if the descended presheaf on $\Hot_\U$ is
representable by a space $Y$, D.8 forces $Y$ to be a $1$-type.
Then $H_K$ is represented by its fundamental groupoid, so the
ind-object is itself essentially constant with a local representative.
\end{proof}

The missing implication is therefore precisely
\[
 \boxed{\quad
 H_K\text{ inverts }\W_\infty
 \quad\Longrightarrow\quad
 H_K\text{ inverts every }\eta_X:X\to\Pi_1X.
 \quad}
 \tag{E.17c}
\]
The reverse implication follows from (E.17a) because
$\W_\infty\subset\W_1$. The maps $\eta_X$ in (E.17c) are
not weak homotopy equivalences when $BX$ has higher homotopy.
Consequently A.4 applied to the genuine groupoid reflection does
not prove (E.17c) from its stated premise. A.7 only proves
$H_K'\cong\colim_i[-,BK_i]$ after descent; the target varies
with $i$, so D.8 cannot be applied to one of the stages as if it
represented that whole descended presheaf.

The size restriction is material. If the index is allowed to be
$\V$-small and unbounded in the $\U$-cardinals, D.7 gives an
invariant ind-presentation of $X\mapsto[BX,K(\Z/2,2)]$.
It cannot be isomorphic to any ind-object of groupoids: every
such ind-object sends the projection from a finite poset of type
$S^2$ to $e$ to a bijection, whereas this functor takes that map
to the nonbijective constant-map function
$\{*\}\to H^2(S^2;\Z/2)=\Z/2$. This is a counterexample to
the version with the enlarged index universe, and does not answer
the $\U$-small or sequential assertions in E.17. Likewise E.4
refutes the later inference about unions even for a diagram already
made of local categories; it does not refute the existence question
posed in E.17.
